\uuid{W0ue}
\titre{Vecteurs}
\auteur{bodin ; hanani ; mzari}
\organisation{exo7}
\licence{CC BY-NC-SA 4.0}
\datecreate{2020-01-20}
\niveau{L1}
\module{Géométrie}
\chapitre{Géométrie du plan}
\sousChapitre{Distance, norme, produit scalaire}
\difficulte{1}
\theme{vecteurs}
\exo7qcmid{lille-1.geometrie-plane.009}

\contenu{
\begin{question}
\qclassification{140.01,140.02}
\qtags{motcle=vecteurs}
Soit $D$ la droite passant par les points $A(1,-1)$ et $B(1,1)$. Quelles sont les bonnes réponses ?
\begin{answers}
    \good{$\vec{u}(0,1)$ est un vecteur directeur de $D$.}
    \bad{$\vec{u}(0,1)$ est un vecteur normal à $D$.}
    \bad{Le point $C(1,0)$ n'appartient pas à $D$.}
    \good{Le point $C(1,0)$ appartient à $D$.}
\end{answers}
\begin{explanations}
Le vecteur $\overrightarrow{AB}=(0,2)$ est un vecteur directeur de $D$. Par ailleurs, $\overrightarrow{AC}=(0,1)=\frac{1}{2}\overrightarrow{AB}$. Donc $C\in D$.
\end{explanations}
\end{question}
}
